\section*{Overview}

The segment tree is the contest default when the problem is really about intervals. It keeps the array in a binary
decomposition of ranges, which makes it flexible enough for sums, minima, maxima, gcds, and many custom merges.

\section*{When to Use It}

Use it when:

\begin{itemize}
  \item queries are on arbitrary intervals,
  \item updates are not limited to prefixes,
  \item the merge operation is associative,
  \item a Fenwick tree is too specialized.
\end{itemize}

It is often the right first structure for path-query frameworks such as heavy-light decomposition too.

\section*{Core Idea}

Every node represents a segment. Its value is the merge of the two child segments. The root stores the whole array,
its children store the left and right halves, and so on until leaves represent single positions.

The reason this is useful is that any query interval can be decomposed into only \(O(\log n)\) disjoint tree nodes.
So instead of scanning the whole interval, we merge a logarithmic number of precomputed summaries.

\section*{Operations}

\begin{itemize}
  \item \textbf{build}: initialize the leaves, then pull parent values upward.
  \item \textbf{point update}: change one leaf, then recompute the ancestors.
  \item \textbf{range query}: walk upward from the left and right borders, collecting the covered segments.
\end{itemize}

The reference implementation uses the iterative layout because it is short, cache-friendly, and hard to break.

\section*{Correctness Intuition}

Two facts do all the work:

\begin{itemize}
  \item every internal node is defined as the merge of its two children,
  \item every query interval can be partitioned into disjoint node intervals from the tree.
\end{itemize}

Once those are true, the answer over any queried interval is just the merge over that partition.

\section*{Complexity Analysis}

\begin{itemize}
  \item Build: \(O(n)\)
  \item Point update: \(O(\log n)\)
  \item Range query: \(O(\log n)\)
  \item Memory: \(O(n)\) in the iterative layout, usually stored as \(2n\)
\end{itemize}

\section*{Implementation}

The code panel shows an iterative sum segment tree with:

\begin{itemize}
  \item build from an array,
  \item point assignment,
  \item point increment,
  \item range sum query on \([l,r]\).
\end{itemize}

If the problem needs a different merge, the structure stays the same and only the node type plus merge function change.

\section*{Common Pitfalls}

\begin{itemize}
  \item Mixing inclusive and half-open interval conventions.
  \item Forgetting that the iterative layout naturally wants 0-indexed positions but stores leaves starting at \texttt{n}.
  \item Using a non-associative merge.
  \item Reaching for a segment tree when the problem is static and a sparse table would be simpler.
\end{itemize}

\section*{Practice Problems}

\begin{itemize}
  \item Dynamic range sum / min queries on an array.
  \item Path queries after heavy-light decomposition.
  \item Segment-tree-based DP transitions where each node stores a custom state.
\end{itemize}

\section*{References}

\begin{itemize}
  \item \href{https://cp-algorithms.com/data_structures/segment_tree.html}{cp-algorithms: Segment Tree}.
  \item \href{https://github.com/kth-competitive-programming/kactl}{KACTL}.
  \item \href{https://usaco.guide/CPH.pdf}{Competitive Programmer's Handbook}.
\end{itemize}
